%!TeX root=../tese.tex
%("dica" para o editor de texto: este arquivo é parte de um documento maior)
% para saber mais: https://tex.stackexchange.com/q/78101

\chapter{Mapeamento Sistemático de impactos ambientais em HPC}

\section{Protocolo e execução}

Enquanto uma revisão sistemática de literatura (RSL) sintetiza evidências a partir da literatura 
para responder a uma questão de pesquisa em específico, um mapeamento sistemático de literatura (MSL) tem como
objetivo produzir uma visão geral e estruturada da pesquisa existente em uma área de interesse. Geralmente, 
mapeamentos sistemáticos tem um escopo mais amplo, sobre uma área de pesquisa significativa, enquanto
revisões sistemáticas tem um escopo mais estreito. Um MSL geralmente resulta no mapeamento e/ou contagem
dos relatos existentes agrupoados em diferentes categorias, enquanto o resultado de uma RSL é 
a síntese do que há de pertinente na literatura acerca daquele questão. O autor desta dissertação
foi um dos coautores do artigo (Lucas), participando nas etapas de triagem e classficação dos relatos
extraídos da literatura. O artigo apresenta um mapeamento sistemático de literatura sobre 
impactos ambientais em HPC, e o presente capítulo é uma adaptação do mesmo.

O MSL foi conduzido de acordo com a metodologia definida por Peterson et al., cobrindo as etapas
de busca, triagem, seleção e classificação dos relatos extraídos da literatura. Os artefatos publicados
tornam possível a reprodução parcial dessas etapas. 

A estratégia de busca foi derivada a partir da lógica \emph{PICO} (\textit{Population, Intervation,
Comparison and Outcome}). \textit{Population} ancora a busca nos sistemas de HPC, enquanto 
\textit{Intervention} captura avaliação ambiental e sustentabilidade. \textit{Comparison} e 
\textit{Outcome} não foram utilizados, pois não se encaixam na natureza de um mapeamento sistemático.
A busca cobriu título, resumo e palavras-chave, e foi realizada em 3 bases de dados: IEEE Xplore,
ACM Digital Library e Scopus. A busca foi realizada com resultados filtrados até 11 de Abril de 2026,
retornando um total de 4812 registros. Após a combinação e remoção de duplicatas, 3788 registros 
foram submetidos à triagem.

Antes da etapa de triagem, se fixaram dois critérios de inclusão e três critérios de exclusão. 
Aqui, o critério mais interessante é o segundo critério de exclusão, que exclui relatos que focam
apenas na questão energética dos sistemas de HPC, sem conectar isso com impactos ambientais
mais amplos. É essse critério que define a fronteira do campo mapeado e que separa o mapeamento
apresentado do resto da literatura já consolidade sobre eficiência puramente energética em HPC.

A triagem combinou três movimentos: primeiro, foi feita uma revisão manual de um subconjunto 
de 365 registros, da qual 25 registros foram considerados para inclusão e então selecionados para o movimento seguinte. Segundo, foram
aplicados passos de \textit{snowballing}, tanto para frente quanto para trás, a partir desses 25 registros. Essas
etapas resultaram na triagem de 2256 registros, dos quais 5 foram retidos. Um terceiro movimento de triagem automatizada foi então
aplicado, utilizando os 365 registros previamente selecionados e os 5 relatos provenientes do \textit{snowballing} como base de 
referência para avalição da triagem automática.

Nessa etapa automatizada, três modelos de linguagem de grande escala (\textit{Large-Language Models}, ou LLMs)
atuaram como avaliadores independentes, onde cada modelo recebia o título, resumo e critérios de inclusão e exclusão,
e devolvia a decisão de inclusão ou exclusão com o critério correspondente. Além disso, os modelos foram instruídos 
a incluir os relatos em caso de insuficiência de informação, para não descartar candidatos que poderiam ser relevantes.
Os três avaliadores automatizados incluíram todos os relatos presentes na base de referência, e o coeficiente de
Kohen contra a mesma variou entre 0,329 a 0,456, o que indica uma concordância levemente moderada entre os avaliadores.

As diferenças entre os avaliadores foram motivadas principalmente por falsos positivos, ou seja, relatos que 
foram incluídos por um avaliador, mas que não atendiam aos critérios de inclusão. Para controlar a quantidade
de falsos positivos, foi definido um critério de inclusão total, onde um relato só seria considerado para análise
posterior caso todos os três avaliadores concordassem com a inclusão. Disso, 514 registros que haviam sido incluídos
por um ou dois avaliadores foram excluídos de revisão humanda, ao passo que 307 registros receberam concordância
de inclusão pelos três avaliadores e seguiram. Desses registros, 244 não pertenciam a base de referência e 
foram revisados conjuntamente pelos dois revisores humanos, com o coeficiente de Kohen de 0,374, indicando uma
concordância moderada entre tais revisores. Após discussões para resolver as divergências, 32 mais relatos foram 
incluídos aos 30 advindos da base de referência, totalizando em 62 relatos. O esquema de classificação desenvolvido
para o mapeamento foi construído a partir da extração das palavras-chave dos resumos e introduções desses 62 relatos,
com as categorias emergindo a partir do agrupamento dessas palavras-chave e sendo posteriormente refinadas
durante a extração. 

\section{Esquema de classificação em quatro facetas}

O processo de extração das palavras-chave ocorreu em várias rodadas, sendo que já no primeiro ciclo foram produzidas
4 facetas que se mostraram consistentes e por isso mantidas pelo restante do processos. As quatro facetas são:
impactos ambientais, estágios do ciclo de vida, locus no sistema, e alavancas de gestão e intervenção. Cada uma
dessas facetas responde a pergunta de pesquisa sobre o relato, sendo elas, respectivamente: qual impacto ambiental
o relato examina, quando no ciclo de vida do sistema de HPC o impacto ocrorre, em qual camada do sistema o relato age,
e como o impacto é tratado ou mitigado pelo relato. 

Três dessas quatro facetas carregam uma ordenção interna, ou seja, formam uma sequência de categorias
ao longo de um eixo com direcionamento. O locus começa pela instalação e infraestruturas de energia,
passa pelo hardware e termina no nível de software e cargas de trabalho. O ciclo de vida vai da produção 
do sistema, passa pela operação do mesmo e termina no estágio de fim de vida. Por fim, as alavancas de gestão
e intervenção progridem da medição para ação, e de mudança técnica para mudança a nível organizacional.

A faceta dos impactos ambientais se divide em quatro categorias, sendo elas: emissões de carbono,
materiais e resíduos, uso de água, e biodiversidade e ecossistemas. Como essa faceta é nominal,
é possível posicioná-la como a faceta central do esquema de classificação, agrupando corpus em populações
e comparando o comportamento de cada uma delas ao longo dos três eixos ordenados das outras facetas.

Esta dissertação está localizada nas categorias de estágio operacional quanto a faceta de ciclo de vida,
no categoria de software e cargas de trabalho quanto a faceta de locus, e na categoria de otimização 
de carga e tempo de execução quanto a faceta de alavancas de gestão e intervenção.

\section{Panorama da literatura}

O corpus não apresenta uma distribuição proporcional de relatos quanto as categorias da faceta de
impactos ambientais, com 54 dos 62 relatos (87,1\%) abordando emissões de carbono,
12 relatos (19,4\%) abordando materiais e resíduos, 5 relatos (8,1\%) abordando uso de água, e
apenas 1 relato (1,6\%) abordando biodiversidade e ecossistemas. A categoria de carbono não só é a mais representada, mas também é a que apresenta a maior parte dos relatos mais recentes, 
passando de 13 relatos publicados entre 2011 e 2022, para 36 relatos publicados apenas entre 2023 e 2025.
Complementando, a maioria desses relatos cobre apenas uma categoria de impacto ambiental, com 52 abordando exatamente um impacto, 10 abordando
dois impactos e nenhum relato abordando três ou mais. Desses 10 relatos que cobrem dois impactos,
9 deles combinam emissões de carbono com materiais e resíduos, enquanto apenas 1 combina emissões
de carbono com o uso de água. 

Esses pares de impactos ambientais que existem no corpus não são frutos de um processo sistemático para captura
de impactos ambientais em HPC, mas sim da existência de conjuntos de dados onde eles coexistem. A avaliação dos 
impactos ambientais de produtos pode ser realizada através da Análise de Ciclo de Vida (em inglês 
\textit{Life Cycle Assessment}, ou LCA), que é uma metodologia padronizada pelas ISO 14040 e 14044.
Esse método contabiliza os impactos ambientais de um produto ao longo de sua existência, passando da
extração de matérias-primas, manufatura, transporte, uso e por fim, descarte. A ACV de manufatura de hardware
com o objetivo de quantificar emissões de carbono acaba por capturar também impactos relacionados a materiais
e resíduos, com essa coexistência dos impactos numa mesma fonte de dados sendo refletida nos 9 relatos em conjunto.
Quanto ao relato que combina emissões de carbono com o uso de água, ele advém de um modelo de resfriamento, que 
captura tanto o consumo de energia (indiretamente, as emissões de carbono) quanto o consumo direto de água para o funcionamento
do sistema de resfriamento.

O estágio operacional é o mais representado para cada uma das quatro categorias de impacto ambiental, 
com abordagem em 52 dos 54 relatos sobre carbono, 10 dos relatos sobre materiais e resíduos, 5 dos 5 relatos 
sobre uso de água, e o único relato sobre biodiversidade. O estágio de produção assume um papel secundário, 
com 15 dos 54 relatos sobre carbono. Essa concentração, assim como no caso sobre a combinação de impactos ambientais,
é reflexo da disponibilidade de dados. Telemetria de potência e traços dos jobs executados nos sistemas de HPC são 
visto como necessários pelos operadores de sistemas, e portanto são produzidos continuamente a partir do momento
que esses entram em produção. Ainda acerca disso, o aumento no número de estudos sobre o estágio de produção é observado,
motivado pelo aumento do número de contabilizações de carbono incorporado. Uma ausência notável quanto a essa
faceta é a falta de relatos sobre o par água-fim de vida. Estudos mostram o impacto que o descarte mal apropriado
de componentes pode causar sobre reservatórios subterrâneos de água, mas nenhum estudo específico sobre
HPC trata disso. Os dados necessários para tais estudos exigiriam uma sinergia com a ecologia industrial, algo
algo ainda não explorado pela comunidade de HPC no que tange o uso de água.

Quanto a faceta de locus, o carbono percorre e tem a maioria dos relatos em todas as categorias dessa faceta:
são 38 relatos situados na camada da instalação, 14 na de hardware e 28 em software e cargas de trabalho. Materiais e resíduos naturalmente
acabam por se concentrar na camada de hardware, com 10 dos 12 relatos situados nessa camada, porém ainda sim possuem 
relatos em outras camadas. Em contraste, o uso de água é abordado apenas na camada da instalação, com todos os 5 relatos 
associados ao resfriamento dos sistemas da instalação. O relato sobre biodiversidade e ecossistemas se encaixa simultaneamente
na camada de hardware e na camada de software. Essas concentrações parecem acompanhar quem produz a pesquisa, 
dado que os relatos sobre materiais permanecem na camada de hardware pois eles são alimentados pelos dados provenientes
dos inventários de ciclos de vida de hardware, enquanto o uso de água permanece atrelado a camada de instalação pois
o trabalho origina-se majoritariamente da engenharia de resfriamento.

A mesma prevalência do carbono é observada para a faceta das alavancas de gestão e intervenção, com 28 relatos em medição e contabilidade,
21 em otimização de carga e tempo de execução, 17 em projeto de infraestrutura e 19 em incentivos e política. Materiais também engajam
nas quatro categorias, dado que estudos de ciclo de vida de hardware geram dados que fluem naturalmente da medição
para possíveis intervenções nas próprias políticas de aquisição de hardware dos centros de HPC. O uso de água engaja duas categorias,
medição em 2 relatos e projetos de infraestrutura em 4 relatos, enquanto o único relato sobre biodiversidade e ecossistemas se encaixa na categoria de incentivos e política.
engaja apenas medição.  

Analisados em conjunto, os três paraeamentos entre impacto ambiental e respectivas facetas descreve um corpus mais
aprofundado ao longo de um corredor único: carbono, estágio operacional, camada da instalação e medição e contabilidade.
Fora dele, principalmente olhando para além das emissões de carbono, cada categoria de impacto ambiental se concentra
de uma forma diferente. Materiais e resíduos se concentram notadamente na camada de hardware, mas conseguem passar por todas
as categorias operacionais e de alavancas, inclusive tendo mais alcance sobre o estágio de fim de vida. Equanto isso,
os relatos sobre uso de água ficam restritos a trabalhos sobre o resfriamento da instalação, sem chegar a hardware ou software.

Os impactos ambientais menos estudados, seja o uso de água ou a biodiversidade, demonstram carecer de instrumentação
comparável entre níveis. Não existem métricas quanto ao consumo de água por um job ou por um nó que corresponda
à telemetria de potência por nó, nem um sinal de rede em tempo real para o uso de água, tão menos um sensoriamento
de biodiversidade em nível de instalação. 

Um escalonamento consciente de água exigiria sinais de consumo hídrico em uma granularidade compatível com os ciclos
de decisão de um escalonador, em uma faixa infra-horária, e é justamente esse canal de dados que não existe. Essa
é a causa principal quanto a lacuna de relatos sobre o uso de água na camada de software e também sobre o uso de água
para alavancas de otimização de carga. A própria medição do consumo de água é tratad como emergente, dado que
o primeiro enquadramento que a trata como pegada de primeira classe em HPC é de 2025. Como referência, relatos sobre carbono
passaram da simples observação para a ação ao redor de uma década, a partir da infraestrutura de sinais de rede
que se tornou disponível no início da década passada. A biodiversidade, por sua vez, enfrenta um desafio ainda maior no momento, visto que há um único ponto de partida de  medição, e
intervenções de otimização ou de governança não podem ser realizadas sem que hajam métricas para quais agir.

É possível inferir portanto que os pontos cegos da leitura são reflexos da fronteira da instrumentação disponível,
não da pegada ambiental de HPC por si. Todos os pares que nenhum relato alcança não correspondem a perguntas 
sem nenhum interesse científico, mas sim a lacuna de dados que a comunidade ainda não construiu.

\section{Limitações do mapeamento}

As limitações do MSL apresentado neste capítulo são organizadas em quatro classes, de acordo com as definições
propostas por Wohlin et al. A primeira limitação é a de construção, que avalia se o mapeamento foi capaz de encontrar
e capturar todos os relatos relevantes sobre impactos ambientiais em HPC. Como a consulta foi aplicada sobre título,
resumo e palavras-chave, se assume que autores de trabalho ambiental expõem esse enquadramento de forma 
explícita nos metadados indexados. Além disso, o esquema de classificação de quatro facetas, que foi construído
de forma indutiva a partir do processo de extração de palavras-chave, é uma projeção defensável mas não a 
única possível de se obter. A segunda limitação é de caráter interno, que avalia a confiabilidade do processo de triagem.
A partir dos coeficientesde Kohen obtidos, tanto entre os avaliadores automatizados quanto entre os revisores humanos,
se assume que a fronteira do corpus foi aproximadamente definida, não de forma exata. A regra dos três votos
deixou 514 registros fora de revisão humana, e o snowballing exluiu 2256 registros apenas a partir do seu título e
resumo, o que favorece relatos cuja faceta de impacto ambiental é mais explícita. A terceira limitação é de caráter
externo, que avalia a capacidade de generalização do mapeamento para além do corpus analisado. Não foram consultadas
Web of Science, DBLP e nem servidores de preprints, o que pode ter deixado de fora relato relevantes. Além disso,
as consultas foram feitas exclusivamente em inglês. Por fim, a quarta limitação é de conclusão, que avalia a
validade das inferências feitas a partir do mapeamento. As contagens são descritivas, sem teste de significância
estatística, e o número enxuto de 62 relatos não sustenta inferência estatística entre os pares de facetas.

Uma limitação importante a ser destacada é de conclusão aplicada aos pares de facetas que não possuem relatos,
como o par de facetas água-software. A ausência de relatos pode indicar tanto uma ausência genuína na literatura,
quanto um problema nas etapas de pesquisa ou extração de palavras-chave. É possível assumir duas premissas que
indicam que a ausência de relatos é genuína, sendo elas: primeiro, um viés na inclusão explicaria quantos relatos
entrariam no corpus, mas não justificaria por que todos que entraram se concentram em um único par de facetas,
como no caso do par água-instalação; segundo, o viés de metadados falha justamente onde não interessa, visto que
um trabalho de escalonamento consciente de água anunciaria água seja no título, resumo ou palavras-chave.

\section{Implicações}

Como mencionado anteriormente, dos 62 relatos que compõe o corpus, apenas 5 deles abordam o impacto do uso de água
causado por centros de HPC, sendo que todos esses se concentram na camada de instalação. Nenhum dos relatos está
situado na camada de software, e nenhum recorre à alavanca de otimização de carga e tempo de execução, cenário que 
contrasta com a situação do carbono, que reúne 28 e 21 relatos respectivamente. Esta dissertação se insere nessa
célula de conchecimento vazio, propondo uma intervenção na camada de software, através da alavanca de otimização
de tempo de execução, atuando sobre o impacto ambiental do uso de água. 

O MSL realizado não apenas localiza a existência dessa célula de conhecimento vazia, mas também fornece uma visão
de por que ela existe. Esse vácuo é reflexo da ausência de instrumentação que permita a medição do uso de água 
em granularida compatível com os ciclos de decisão de um escalonador. A partir disso, é possível inferir 
a form de contribuição, visto que antes de qualquer intervenção, é necessário construir a instrumentação 
que as torna avaliáveis. Por isso este texto apresenta em primeiro lugar o modelo de pegada e ambiente de testes,
e em seguida as heurísticas que intervém a partir dessa instrumentação construída. 

Como apenas um dos relatos de todo o corpus pareia carbono com água, contabilizar e intervir as duas pegadas simultaneamente
e com preocupações separadas se mostra mais do que um refinamento incremental, mas sim uma necessidade da área. O corpus demonstrou
a combinação de dois impactos ambientais em alguns relatos, a partir da lógica de coexistência dos dados. A partir dessa mesma
lógica, é possivel obter as intensidades de carbono e de água a partir da composição instantânea da rede elétrica
em conjunto com os fatores de emissão de carbono e de consumo hídrico das diferentes tecnologias de geração de energia
elétrica.  Por fim, a justificativa quanto a reportar as pegadas de carbono e de água separadamente segue a mesma 
lógica do relato que combina os dois impactos. Vários sistemas de resfriamento visam reduzir suas emissões de carbono,
e fazem isso através da redução do consumo de energia elétrica. Porém, para alcançar essa redução, boa parte desses sistemas
acaba necessitando de um volume considerativo de água. Coincidentemente, as intensidades de carbono e de água da rede
elétrica podem variar em direções opostas, de forma que essa discordância é um achado importante para a comunidade,
não apenas um ruído de medição.

O mapeamento sugere que ganhos de eficiência energética, por si só, são repetidamente absorvidos pelo próprio
crescimento initerrupto das demandas computacionais, o que sugere que a eficiência energética isolada não é suficiente
para se alcançar um crescimento sustentável. Apesar desta dissertação se situar como uma intervenção de eficiência,
a pergunta primária é sobre demarcação, quanto existe de espaço para economizar, e que fração desse espaço
é alcançável através de políticas simples. Além disso, ao estabelecer que a energia ativa é invariante ao escalonamento
sob premissas declaradas, um limite é estabelecido, cujo limite serve igualmente a quem pretende utilizar a
alavanca de otimização de carga e também a quem pretende demonstrar que ela é insuficiente. Por fim, a contribuição
da instrumentação preenche a célula de medição de uso de água na camada de software e cargas de trabalho.
Consequentemente, o produto deste capítulo não é apenas a constatação de uma lacuna, mas sim o entendimento do que
causa essa lacuna, entendido aqui como a ausência de instrumentação, e a proposição de que tal instrumentação é construível.